\section*{Author power-series and complexity conventions}
A formal power series is a generalization of a polynomial, where the number of terms can be infinite. Formally, $A=\sum_{i\ge0}A_ix^i$ with $A_i\in\mathbb F$, is a power series in the power series ring $\mathbb F[[x]]$. We define the degree $d$ truncation $\mathrm{trunc}(A,d):=\sum_{0\le i\le d}A_ix^i$.
\paragraph{Definition 12 (Algebraic and Transcendental Power Series).} A formal power series $f\in\mathbb C[[x]]$ is said to be algebraic if there exists a polynomial $g\in\mathbb C[x][t]$ such that $g(f)=0$. Otherwise $f$ is said to be transcendental.
For a polynomial $f\in\mathbb F[\mathbf x]$, $L_{\mathbb F}(f)$ is the length of the smallest division-free arithmetic circuit (with only $\{+,-,\times\}$ gates) computing $f$. Constant-free complexity $\tau(f)$ is the size of the minimal constant-free circuit computing $f$, with constants in $\{-1,0,1\}$. Straight-line programs for integers start at 1 and use $+,-,\times$; $\tau(N)$ is the minimum such length for $N$. Subscripts are omitted when the field is clear. These are the author conventions from Sections2/5.2.
\paragraph{Definition 36 (The Stern sequence).} The sequence $(a_n)_{n\ge0}$ given by $a_0=0,a_1=1$, and when $n\ge1$, by $a_{2n}=a_n$ and $a_{2n+1}=a_n+a_{n+1}$, is called the Stern sequence. Its generating function is $A(x):=\sum a_nx^n$.
\paragraph{Theorem 37 (Lemma2.1 and Theorem2.2 in [12]).} If $A(z)$ is the generating function of the Stern sequence, then (1) $A(x^2)=A(x)\frac{x}{x^2+x+1}$; (2) the function $A(x)$ is transcendental.
\paragraph{Multiplicative sequence.} The sequence $(f_n)_{n\ge0}$ is defined as: $f_0=1,f_1=1,f_2=-1,f_p=1$ for all odd primes $p$ and $f_{ab}=f_af_b$. We look at the corresponding generating function $F(x):=\sum f_nx^n$.
\paragraph{Theorem 39 ([13,Theorem2]).} The power series $F(x)$ is transcendental.
\paragraph{Holonomic sequence.} A sequence $(h_n)_{n\ge0}$ is called holonomic if it satisfies a recurrence of the form $a_r(n)h_{n+r}+a_{r-1}(n)h_{n+r-1}+\cdots+a_0(n)h_n=0$, where $a_i$ are polynomials in $n$. The corresponding generating function $H(x):=\sum h_nx^n$ is said to be a holonomic function.
\paragraph{Lemma 42 ([28]).} If $F(x)=\sum_{n\ge0}f_nx^n$ is a power series in $\mathbb C[[x]]$ and the radius of convergence of $F(x)$ is zero then $F(x)$ is transcendental. This follows directly from Proposition2 in [28], as stated by the authors.
